% ============================================================
% f2013_fujian_science_q18_parabola.tex
% 2013 福建卷（理工类）第 18 题 正方形与抛物线构造（TikZ 矢量重绘）
% O(0,0), A(10,0), B(10,10), C(0,10)
% OA、AB 各十等分；A_i(i,0)，B_i(10,i)；OB_i 与过 A_i 的竖线交于 P_i
% P_i(i, i^2/10) 在抛物线 y=x^2/10 上
% ============================================================
\begin{tikzpicture}[
  x=0.42cm, y=0.42cm, line width=0.6pt, font=\footnotesize, >=Stealth
]

  % ---------- 正方形 ----------
  \draw (0,0) -- (10,0) -- (10,10) -- (0,10) -- cycle;

  % ---------- 十等分刻度 ----------
  \foreach \i in {1,...,9} {
    \draw (\i,0) -- (\i,-0.22);
    \draw (10,\i) -- (10.22,\i);
  }

  % ---------- OB_i 与竖线 ----------
  \foreach \i in {1,...,9} {
    \draw[gray!65, line width=0.35pt] (0,0) -- (10,\i);
    \draw[dashed, gray!65, line width=0.35pt] (\i,0) -- (\i,{\i*\i/10});
    \fill (\i,{\i*\i/10}) circle (0.09);
  }

  % ---------- 标注 ----------
  \node[below left]  at (0,0)   {$O$};
  \node[below right] at (10,0)  {$A$};
  \node[above right] at (10,10) {$B$};
  \node[above left]  at (0,10)  {$C$};

  \node[below] at (1,-0.28) {$A_1$};
  \node[below] at (2,-0.28) {$A_2$};
  \node[below] at (5,-0.28) {$A_i$};
  \node[below] at (9,-0.28) {$A_9$};

  \node[right] at (10.28,1) {$B_1$};
  \node[right] at (10.28,2) {$B_2$};
  \node[right] at (10.28,5) {$B_i$};
  \node[right] at (10.28,9) {$B_9$};

  \node[right] at (5.35,2.95) {$P_i$};

\end{tikzpicture}
